\section{Methodology}

\subsection{Problem Formalization}

Mathematical formalization of agent-induced pricing distortions. Formal definition of potential loss mechanisms $\alpha D$

We consider a business across time during which we have an evolving vector $p_t \in \Re^N$ where $N$ is the number of products in our catalogue. our price vector is directly dependent on a demand function $q_t$ which we define as a linear method of a price elasticity matrix $B_t$. This is the same setup that Microsoft created in their research.

We gether interaction data from users interacting with a sample platform simulating a hotel/airline which generates interaction distributions $I_t = \{(p_t, q_t^\text{obs}, \pi_t)\}_{t=1}^T$


\subsection{Cost of Information Framework}

Mathematical demonstration and validation of the COI and citation backed evidence, and framework overview + show harm to user via other cost distortions. Maybe split into 3.2.1 (COI Theory) and 3.2.2 (Framework Design)

\subsection{System Architecture}

In order for our research to have grounding in interactions we built a robust e-commerce web-platform. We initially conducted a survey of the leading platforms of airlines and hotel booking sites to identify the specific interface patterns that effectively manage complex travel data. Our analysis revealed a clear industry standard: while both sectors rely on tabbed service selection and left-sidebar filtering to streamline navigation, they diverge in result presentation—airlines utilize visual date-price bars and multi-step wizards to optimize for logistical transparency, whereas hotel platforms leverage image-led cards and scarcity triggers to drive emotional engagement and urgency. Our web framework defines a highly agnostic boilerplane which can be seeded with any data-modality with an easy-to-tailor pattern, which we leverage to define a \texttt{hotel} and \texttt{airline} mode. Both modes are then individually deployed via an envrionment level argument which adjusts the proxy routing with a custom middleware inside next.js to render only the desired mode. The purpose of this was to create a baseline adaptable to any use-case or desired commercial application.




\subsection{Experimental Design}

\begin{figure}[ht]
  \resizebox{\columnwidth}{!}{%
    \definecolor{mygreenfill}{RGB}{169, 234, 186}
\definecolor{mygreenborder}{RGB}{29, 145, 61}
\definecolor{mybluefill}{RGB}{204, 222, 255}
\definecolor{myblueborder}{RGB}{66, 106, 189}
\definecolor{mygray}{RGB}{150, 150, 150}



\begin{tikzpicture}[
    node distance=2cm,
    % Style for Green Nodes
    greenbox/.style={
        rectangle,
        draw=mygreenborder,
        fill=mygreenfill,
        line width=1.2pt,
        align=center,
        minimum height=1cm
    },
    % Style for Blue Nodes
    bluebox/.style={
        rectangle,
        draw=myblueborder,
        fill=mybluefill,
        line width=1.2pt,
        align=center,
        minimum height=1cm
    },
    % Style for Arrows
    myarrow/.style={
        ->,
        >={Stealth[length=3mm, width=2mm]},
        draw=black!80,
        line width=1.2pt,
        rounded corners=5pt
    },
    % Style for Background Dashed Circles
    dashedloop/.style={
        dashed,
        draw=mygray,
        line width=1pt
    }
]

    % --- Coordinate Layout ---
    % Defining a grid relative to the center

    % Left Loop (Green) Nodes
    \node[greenbox, minimum width=3.5cm] (commerce) at (-3.5, 2) {Commerce Experiment};
    \node[greenbox, minimum width=1.5cm] (raw) at (-6.5, 0) {Raw\\Logs};
    \node[greenbox, minimum width=1.5cm] (features) at (-4, -2.5) {Features};
    \node[greenbox, minimum width=2.5cm] (classification) at (-1, -0.5) {Classification\\Training A/H};

    % Right Loop (Blue) Nodes
    \node[bluebox, minimum width=2.5cm] (trainedpricing) at (3.2, 2) {Trained Pricing};
    \node[bluebox, minimum width=2.5cm] (policy) at (6.5, 0) {Trained Pricing\\Policy};
    \node[bluebox, minimum width=2.5cm] (rlgym) at (3.2, -2.2) {RL Gym\\Training};

    % --- Background Dashed Loops ---
    \begin{scope}[on background layer]
        % Left Loop Circle
        \draw[dashedloop] (-3.5, 0) ellipse (3.5cm and 2.8cm);
        % Right Loop Circle
        \draw[dashedloop] (3.5, 0) ellipse (3.5cm and 2.8cm);
    \end{scope}

    % --- Arrows: Loop One (Green) ---
    % Commerce -> Raw Logs
    \draw[myarrow] (commerce.west) to[out=180, in=90] (raw.north);

    % Raw Logs -> Features
    \draw[myarrow] (raw.south) to[out=270, in=180] (features.west);

    % Features -> Classification
    \draw[myarrow] (features.east) to[out=0, in=250] (classification.south);

    % Classification -> Commerce (Closing the loop)
    \draw[myarrow] (classification.north) to[out=110, in=0] (commerce.east);

    % --- Arrows: Loop Two (Blue) ---
    % Classification (Green) -> RL Gym (Blue) - Crossing over
    \draw[myarrow] (classification.east) to[out=0, in=180] (rlgym.west);

    % RL Gym -> Policy
    \draw[myarrow] (rlgym.east) to[out=0, in=270] (policy.south);

    % Policy -> Trained Pricing
    \draw[myarrow] (policy.north) to[out=90, in=0] (trainedpricing.east);

    % Trained Pricing -> Commerce (Crossing back)
    \draw[myarrow] (trainedpricing.west) -- node[above, font=\small, yshift=2pt] {New Pricing} (commerce.east);

    % --- Text Labels ---

    % Loop One Label
    \node[align=center] at (-3.8, 0) {Loop One:\\Data \textit{(Online)}};

    % Loop Two Label
    \node[align=center] at (3.5, 0) {Loop Two:\\Defense Gym \textit{(Offline)}};

    % Bottom Legend
    \node[font=\small] (taskA) at (-4, -4) {Dynamic Pricing Task A};
    \node[font=\small] (taskB) at (4, -4) {Dynamic Pricing Task B};
    \node[font=\small] (indep) at (0, -4) {Independent};

    % Arrows for bottom legend
    \draw[->, >=Stealth, thick, darkgray] (indep.west) -- (taskA.east);
    \draw[->, >=Stealth, thick, darkgray] (indep.east) -- (taskB.west);

\end{tikzpicture}

  }
  \caption{Overview of the Dynamic Pricing Tasks.}
\end{figure}
Study methodology and approach. Data acquisition strategy. Defined objectives and success criteria. Observable metrics and KPIs

\subsection{Dynamic Pricing Algorithm Analysis}
Deep dive into how the algorithm works, different kinds and justification for chosen appraoches + agent impact modeling and quantification.
\subsection{Reinforcement Learning Formulation}
How do we define the state space, action space and reward function breakdown and algorithm benchmarking.
POSSIBLY: Expand into full subsections: 3.6.1 (State-Action Space), 3.6.2 (Reward Design), 3.6.3 (Benchmarking)


\begin{algorithm}[t]
\DontPrintSemicolon
\KwIn{stepsize $\eta$, smoothing $\delta$, rank $d$}
\For{$t=1$ \KwTo $T$}{
  Sample $u_t$ on unit sphere; set $x_t^\prime=x_t+\delta u_t$\;
  Set $p_t \gets U x_t^\prime$ and observe $q_t, R_t(p_t)$\;
  $x_{t+1} \gets \Pi\_{\mathcal{X}}(x_t-\eta R_t(p_t) u_t)$\;
}
\caption{Online Pricing Optimization (template)}
\end{algorithm}
